\section*{Abstract}

Modern automated theorem provers are very effective, but whether their results can be trusted depends on an independent check of the proof objects they produce.
This thesis presents the design and implementation of a Python-based proof verifier for TPTP/TSTP derivations, developed according to the specifications of ProoVer 2026, a proof-checking competition held as part of the FLoC Olympics 2026 in Lisbon~\citep{proover-competition}.

In proof verification, a wrongly accepted proof is by far the worst outcome, and the scoring scheme of the competition reflects this: a wrong acceptance costs ten times as much as a correct acceptance gains~\citep{proover-competition}.
The ruleset of ProoVer 2026 also prescribes a three-valued verdict, with \emph{VerifiedGood} for a valid proof, \emph{VerifiedBad} for an invalid proof, and \emph{Unknown} if the checker could not reach a definite conclusion.
A fourth answer, \emph{Timeout}, is given if the verifier does not finish within the time limit of 30 seconds.
The verifier developed in this thesis accepts a proof only if every step has been validated, rejects it only if it has found a concrete defect, and answers \emph{Unknown} in all other cases.
This policy is implemented as a modular pipeline of parsing and normalization, structural and dependency checks, and rule-specific inference checks.
Structural, propositional, quantifier, and clausal rules as well as Skolemization are checked internally, and every failing step gets its own diagnostic.
All remaining steps are passed to an external prover, in this thesis, Vampire, which only sees the premises of the step and never the proof as a whole.

Of the \EvalOffN{} problem/proof pairs provided by the competition committee, the verifier accepts all \EvalOffValid{} valid proofs, rejects \EvalOffRejInvalid{} of the \EvalOffInvalid{} invalid proofs, and answers \emph{Unknown} for the remaining one.
The median run time of \EvalWallMedian{} per proof is well within the limit of 30 seconds wall clock time~\citep{proover-competition}.
This gives \EvalOffScore{} of \EvalOffMaxScore{} points, and no proof is wrongly accepted.
The same holds for the \EvalSampleN{} proofs of the originally published set of 100 that were not scored in the competition.
A generated suite of \EvalExtraN{} rule-level test proofs and \EvalMutN{} mutants of valid proofs, in which every mutant that is no longer valid is rejected, suggests that the checks are not simply tuned to the benchmark.
Even when the external prover is removed, the verifier accepts no invalid proof.
The contributions of this thesis are a soundness argument, a modular verification architecture, a working implementation of the verifier, and an empirical analysis of its performance and robustness under simulated competition conditions.

\textbf{Keywords}: Automated Theorem Proving, Proof Verification, TPTP, TSTP, Fail-Closed Checking, Skolemization, ProoVer 2026.
